% _extensions/gb9704/context-template.tex
% 字体方案（必须在 \usemodule 之前设置）：
%   reference  — GB/T 9704 参考文档（SimHei/楷体国标/仿宋国标）
%   fangzheng  — 方正全套（方正黑体/方正楷体/朱雀仿宋）
%   modern     — 现代简洁（文泉等宽正黑/楷体国标/朱雀仿宋）
$if(gbtfontpreset)$\def\gbtfontpreset{$gbtfontpreset$}$else$\def\gbtfontpreset{modern}$endif$

% 加载 gbt9704 模块（字体/CJK断行/标题样式已由模块统一管理）
\usemodule[gbt9704]

% ==========================================================
% 补充字体定义 — 与 LaTeX gbt9704.cls 字体一致
% 将黑体注册到 mainface 的 ss（sans-serif）通道。
% ⚠️ 不调用 \setupbodyfont 刷新——会破坏模块的 \definefont 独立字体。
% ConTeXt LMTX 存在字体 ID 污染 bug：\definefont 字体与 fontfamily
% 通道互斥，先使用者污染后者。section 标题优先级最高，通过 \ss 通道保障。
% ==========================================================
\definefontfamily[mainface][ss][simhei][features=default]
\definefont[SimSun][name:simsun*default]

% ==========================================================
% emoji 支持 — Pandoc Lua filter 调用 \emoji{} 时使用
% ==========================================================
\definefont[EmojiFont][name:notocoloremoji*default]
\def\emoji#1{{\EmojiFont #1}}

% ==========================================================
% 修复 \redseparator — 模块用了 LaTeX 语法，ConTeXt 不支持
% ==========================================================
\def\redseparator{%
  \blank[2pt]
  \noindent\blackrule[color=officialred, width=\textwidth, height=1.5pt]
  \blank[4pt]
}

% ==========================================================
% 覆盖标题样式 — 与 LaTeX gbt9704.cls 一致
% GB/T 9704: 一级黑体、二级楷体、三级仿宋加粗
% section 用 \ss（mainface ss 通道→simhei/黑体），绕过模块的
% \definefont[SimHei]，确保字体嵌入 PDF（否则 ConTeXt LMTX 字体 ID 污染）
% ==========================================================
\setuphead[title][
  style={\XiaoBiaoSong\switchtobodyfont[22pt]},
  color=officialred,
  align=middle,
  number=no,
]
\setuphead[section][
  style=\ss,
  number=no,
  align=flushleft,
  indentnext=yes,
  margin=2em,
]
\definefont[GBKaiBold][name:kaitigb2312*fakebold]
\setuphead[subsection][
  style={\GBKaiBold},
  number=no,
  indentnext=yes,
  margin=2em,
]
\setuphead[subsubsection][
  style={\FakeBold},
  number=no,
  indentnext=yes,
  margin=2em,
]

% ==========================================================
% 列表环境别名（兼容 Pandoc 默认语法）
% ==========================================================
\let\startenumerate\startitemize
\let\stopenumerate\stopitemize
\definedescription[description]

% ==========================================================
% 文本装饰命令（docx2ctx 生成器使用 \bold, \emph, \scaps, \overstrike, \high, \low）
% gbt9704 模块锁定命令名，用 \pushoverloadmode 允许覆盖
% ==========================================================
\pushoverloadmode
\def\bold#1{\FakeBold{#1}}
\def\emph#1{\FakeSlant{#1}}
\def\scaps#1{{\sc#1}}
\popoverloadmode

% ==========================================================
% 三线表样式（覆盖 gbt9704 模块默认）
% ==========================================================
\setupTABLE[
  frame=off,
  rulethickness=0.5pt,
  align={flushleft,lohi},
]
\setupTABLE[r][first][topframe=on, bottomframe=on]
\setupTABLE[r][last][bottomframe=on]
\setupTABLE[c][1][background=]     % 清除模块的灰色表头背景

% ==========================================================
% 标注块（Callouts）环境 — 对齐 Quarto LaTeX tcolorbox 样式
% ==========================================================
\definecolor[callout-note-bg]    [r=0.89, g=0.93, b=0.98]
\definecolor[callout-note-fg]    [r=0.27, g=0.51, b=0.93]
\definecolor[callout-warning-bg] [r=0.98, g=0.93, b=0.82]
\definecolor[callout-warning-fg] [r=0.94, g=0.68, b=0.31]
\definecolor[callout-tip-bg]     [r=0.80, g=0.96, b=0.88]
\definecolor[callout-tip-fg]     [r=0.01, g=0.72, b=0.46]

% 标注块：白底 + 彩色细框（简洁风格）
\defineframedtext[gbtnote][frame=on,framecolor=callout-note-fg,rulethickness=0.4pt,offset=6pt,background=color,backgroundcolor=white,width=local,before={\blank[line]},after={\blank[line]}]
\defineframedtext[gbtwarning][frame=on,framecolor=callout-warning-fg,rulethickness=0.4pt,offset=6pt,background=color,backgroundcolor=white,width=local,before={\blank[line]},after={\blank[line]}]
\defineframedtext[gbttip][frame=on,framecolor=callout-tip-fg,rulethickness=0.4pt,offset=6pt,background=color,backgroundcolor=white,width=local,before={\blank[line]},after={\blank[line]}]
\defineframedtext[gbtimportant][frame=on,framecolor=darkred,rulethickness=0.4pt,background=color,backgroundcolor=white,corner=round,radius=1mm,width=local,before={\blank[line]},after={\blank[line]}]
\defineframedtext[gbtcaution][frame=on,framecolor=darkorange,rulethickness=0.4pt,background=color,backgroundcolor=white,corner=round,radius=1mm,width=local,before={\blank[line]},after={\blank[line]}]

% 标注块别名（兼容 notes 模板命名）
\let\startnoteblock\startgbtnote  \let\stopnoteblock\stopgbtnote
\let\startwarningblock\startgbtwarning  \let\stopwarningblock\stopgbtwarning
\let\starttipblock\startgbttip  \let\stoptipblock\stopgbttip
\let\startimportantblock\startgbtimportant  \let\stopimportantblock\stopgbtimportant
\let\startcautionblock\startgbtcaution  \let\stopcautionblock\stopgbtcaution

% ==========================================================
% 文档开始
% ==========================================================

\starttext

% ⚠️ 预初始化 \ss（黑体）—— ConTeXt LMTX 字体 ID 污染 workaround：
% 必须在所有其他字体之前使用，确保黑体在 section 标题中嵌入 PDF。
% 若先使用其他 \definefont 字体（如 \DaBiaoSong），\ss 将被污染失效。
{\ss\relax}

$body$

\stoptext
